\section{Introduction}
\label{sec:introduction}

The global transition toward environmentally sustainable refrigeration, heat pumps, and industrial thermal management represents one of the most pressing decarbonization challenges in modern chemical engineering. Under the Kigali Amendment to the Montreal Protocol and associated international climate mandates, traditional hydrofluorocarbon (HFC) working fluids---which exhibit atmospheric lifetimes spanning decades and global warming potentials (GWPs) thousands of times that of $\text{CO}_2$---are undergoing aggressive regulatory phasedowns. Consequently, next-generation, ultra-low-GWP hydrofluoroolefins (HFOs) and hydrochlorofluoroolefins (HCFOs) are rapidly emerging as drop-in replacements. Simultaneously, the circular reclamation of spent fluorinated gases, the extractive separation of near-azeotropic refrigerant mixtures, and the development of energy-efficient absorption refrigeration cycles necessitate the deployment of non-volatile, thermodynamically stable solvent media. Ionic liquids (ILs), characterized by negligible vapor pressures, broad liquid operating ranges, and synthetically tunable cation--anion architectures, have garnered intense interest as tailored green absorption solvents and entrainers~\cite{Qin2024_CAILD}.

However, the molecular design space for ionic liquid-mediated refrigerant absorption systems is vast and complex. Pairing thousands of synthetically accessible IL cation--anion combinations with emerging fluorinated refrigerant solutes across continuous operating temperature ($T$) and pressure ($P$) envelopes generates a combinatorial space that defies exhaustive experimental characterization. Classical thermodynamic modeling frameworks offer incomplete solutions: standard group-contribution equations of state (such as UNIFAC) frequently lack interaction parameters for newly synthesized fluorinated groups and asymmetric ionic functional moieties~\cite{Fredenslund1975_UNIFAC,Lei2014_UNIFAC_IL}, while quantum-chemical conductor-like screening models (such as COSMO-RS) incur substantial computational expense per mixture and exhibit known quantitative transfer errors when predicting complex hydrogen bonding and halogen-mediated dispersion interactions~\cite{Klamt2000_COSMORS}.

In response, data-driven machine learning---and particularly molecular Graph Neural Networks (GNNs)---has emerged as a transformative paradigm for chemical property prediction. By parameterizing molecular entities as topological graphs and iteratively updating atomic representations via message-passing algorithms, modern deep GNNs can learn complex non-linear structure--property mappings directly from molecular graphs without requiring manual descriptor engineering. In a prominent recent study, Chu et al.~\cite{Chu2024_AEGNN} demonstrated the utility of Attention-Enhanced Graph Neural Networks (AEGNN) operating on disconnected cation, anion, and refrigerant graphs under ambient operating conditions, reporting remarkable interpolation accuracy ($R^2 > 0.98$) on randomly partitioned HFC/HFO--IL solubility benchmarks.

\textbf{Yet, a critical methodological crisis underlies current molecular AI paradigms in chemical thermodynamics:}

\begin{enumerate}
    \item \textbf{The Fallacy of Random Splits in Multicomponent Systems}: Standard benchmark evaluations routinely rely on uniformly random train/test partitioning. In ternary fluid phase equilibria ($\text{cation} + \text{anion} + \text{gas}$), random splitting allows structural fragments of test solutes, cations, and anions to appear across both partitions, while repeated measurements of identical mixtures across temperature/pressure grids leak into the training corpus. As demonstrated by Li et al.~\cite{Li2023_NatCommun} in ionic polymer electrolytes and Choudhary et al.~\cite{Choudhary2024_OODBenchmark} in materials property prediction, compositional and structural redundancy under random splitting artificially inflates predictive metrics and masks catastrophic failure under genuine Out-of-Distribution (OOD) chemical shifts~\cite{Wu2018_MoleculeNet}.
    
    \item \textbf{The Collapse of the Single-Scalar Distance Assumption}: A foundational premise of modern cheminformatics is the ``applicability domain'' hypothesis, which posits that generalization error increases monotonically with distance from the training distribution in some scalar metric space, such as Tanimoto fingerprint similarity ($D_{\rm FP}$) or Euclidean descriptor distance ($D_{\rm phys}$)~\cite{Tropsha2003_QsarValidation,Netzeva2005_ApplicabilityDomain}. However, \textbf{our systematic mapping of the continuous generalization landscape reveals that no single scalar distance provides a sufficient description of prediction difficulty across multi-dimensional chemical distribution shifts}. Solutes situated at severe topological separation can be predicted with high fidelity, whereas structural isomers with near-zero scalar distance can suffer severe predictive degeneration. Standard scalar metrics are blind to discrete quantum-mechanical dipole cancellations (e.g., symmetric R134 vs. asymmetric R134a) and terminal fluorination shifts that govern liquid-phase phase equilibrium.
    
    \item \textbf{The Unresolved Role of Thermodynamic Grounding}: While augmenting neural networks with physical descriptors has gained traction~\cite{Karpatne2017_PGML,Karniadakis2021_PINN_Review}, physical priors are frequently applied in an unprincipled manner. It remains fundamentally unresolved whether physics regularization acts as a universal inductive regularizer, or whether its benefit is strictly selective depending on the dominant failure mode of the distribution shift.
    
    \item \textbf{The Unseen Representation-Support Boundary}: Standard 2D molecular graph featurization completely omits geometric isomerism, rendering cis/trans isomers mathematically degenerate ($\mathcal{G}_E \equiv \mathcal{G}_Z$). While 3D equivariant architectures can theoretically resolve stereochemical equivalence~\cite{Morehead2024_GCPNet,Adams2024_GeomComplete,Schwaller2024_3DGeom}, the fundamental interplay between architectural expressivity and training data distribution support remains uncharacterized in chemical property prediction.
\end{enumerate}

\subsection*{Contributions of this Work}
In this study, we address these foundational challenges through an exhaustive, diagnostically grounded investigation across $4,444$ high-precision vapor--liquid equilibrium measurements. Our contributions are fivefold:

\begin{itemize}
    \item \textbf{The Chemical Generalization Ladder}: We formulate a five-tier benchmark ladder probing orthogonal distribution shifts across solute identity (M1 LORO, $12$ pure HFCs), anion structural classes (B1 fluorosulfonates and B2 inorganic fluorides), compositional recombination (L2 unseen ionic pairs), and cross-chemical-family saturation transfer (HFO/HCFO zero-shot probe, $N=1,106$).
    
    \item \textbf{The Selective Thermodynamic Grounding Principle}: Grounding regression models in the van der Waals Theorem of Corresponding States~\cite{Pitzer1955_Acentric,Prausnitz1998_FluidPhase,Kocis2020_PINN_Thermo}, we formulate a reduced thermodynamic model ($M_{\rm reduced}$) parameterized by reduced temperature ($T_r = T/T_c$), reduced pressure ($P_r = P/P_c$), and the Pitzer acentric factor ($\omega$). We demonstrate that thermodynamic coordinates do not uniformly boost accuracy; rather, their benefit is \textbf{strictly selective}, slashing error by $36.5\%$ on solute extrapolation and $30.4\%$ on unsaturated zero-shot transfer, while exerting negligible effect ($-0.7\%$) when topological graph redundancy is already dense.
    
    \item \textbf{Mechanistic Diagnostic Attribution of Representation Bottlenecks}: Using Integrated Gradients attribution, continuous landscape mapping, and quantum-chemical GFN2-xTB calculations, we reveal that generalization failures are driven by discrete structural transitions---such as internal dipole moment cancellation in R134 versus R134a ($M_{\rm reduced}\text{ MAE}: 0.1414\text{ vs } 0.0769$) and backbone fluorination asymmetry in R245fa versus R236fa---that isotropic distance metrics fail to capture.
    
    \item \textbf{Automated Epistemic Uncertainty Guardrails}: We demonstrate that deep ensemble disagreement across independently trained seeds serves as an actionable epistemic-uncertainty proxy~\cite{Lakshminarayanan2017_DeepEnsembles}, exhibiting consistent rank correlation with true error ($\rho \approx 0.31\text{--}0.61$) and enabling selective prediction regimes~\cite{Geifman2017_SelectiveNet} that reduce operating risk by over $50\%$ at 50% coverage.
    
    \item \textbf{The Representation-Support Boundary Audit}: We formalize the 2D graph stereochemical degeneracy problem on R1336mzz(E/Z) isomers and execute a controlled intervention experiment, establishing that architectural expressivity cannot overcome the absence of training distribution support.
\end{itemize}

Together, these findings shift the paradigm from empirical benchmark chasing to diagnostic, physics-grounded representation learning for chemical engineering thermodynamics.
